\section{Background}

We briefly provide notions about %Large Language Models, 
natural language ambiguities, Concurrent Game Structures, ATL and ATL$^*$ (analyzing both syntax and semantics). We assume some familiarity with standard concepts from temporal logic and multi-agent systems, focusing on the essential aspects for our framework.


\paragraph{Natural Language Ambiguities.}
Linguistics Formal Syntax and Semantics offer well defined formalizations for identifying natural language ambiguity drawing on graph theory and logical analysis \cite{chomsky1970remarks}. We focus on syntactic ambiguities in translation, as they can be reliably detected and verified. Semantic ambiguities arising from subtle differences in the meaning of terms or expressions present a greater challenge. They appear in the dataset but determining the correct interpretation is non trivial. Consequently, accuracy metrics for syntactic ambiguities are more reliable than those for semantic ones (see the Evaluation section for more details). Here we address the ambiguities covered in our work. The first two types are  right and left dislocation, which occurs when a constituent, wether an argument or adjunct appears outside the clause boundaries, either preceding or following it. Right dislocation is illustrated by "They can eventually ensure the game ends, the players", while "The vending machine, Mary can ensure it will work" exemplifies left dislocation. 
\begin{figure}[h]
\centering
\scriptsize
\begin{forest}
for tree={
    grow'=south,
    parent anchor=south,
    child anchor=north,
    align=center,
    l sep=1pt,   % distanza verticale
    s sep=1pt     % distanza orizzontale
}
[S
  [S
    [NP [They]]
    [VP
      [V [can\\eventually\\ensure]]
      [CP [the game\\ends]]
    ]
  ]
  [{NP$_{\mathrm{disl}}$}
    [the players]
  ]
]
\end{forest}
\caption{Right Dislocation. Cartographic analysis provided in Appendix A.}
\end{figure}

\begin{figure}[h]
\centering
\scriptsize
\begin{forest}
for tree={
    grow'=south,
    parent anchor=south,
    child anchor=north,
    align=center,
    l sep=10pt,
    s sep=8pt
}
[S
  [{NP$_{\mathrm{disl}}$}
    [The vending\\machine]
  ]
  [S
    [NP [Mary]]
    [VP
      [V [can\\ensure]]
      [CP [it will\\work]]
    ]
  ]
]
\end{forest}
\caption{Left Dislocation. Cartographic analysis provided in Appendix A.}
\end{figure}
The third type is Verb Phrase ellipsis (VP ellipsis), a phenomenon whereby a verb phrase is left out of a sentence when its meaning can be recovered from the phrasal context. This typically happens when omitted material mirrors a verb phrase already introduced earlier in the discourse; an example follows "Amanda has a strategy to surely sell all the toys and Leo does too". The fourth type is Right Node Raising (RNR), a construction in which a shared argument appears at the right periphery of a coordinate structure, as in "Bob has, but Albert has not, a strategy to ensure that the vending machine is always operative". The fifth and final type is Quantifier Scope Ambiguity (QSA), where a sentence like "Every agent can prevent a breach" admits two readings. The first is "For every agent, there is some breach that the agent can prevent (possibly a different breach for every agent)". The second is "There is one breach such that every agent can prevent that breach". The respective logical forms in First Order Logic follows:
\begin{itemize}
    \item[-] First reading: $\forall x \, \bigl(A(x) \rightarrow \exists y \, (B(y) \land P(x, y))\bigr)$
    \item[-] Second reading: $\exists y \, \bigl(B(y) \land \forall x \, (A(x) \rightarrow P(x, y))\bigr)$
\end{itemize}
with $A$ being the monadic predicate "is an agent", $B$ the monadic predicate "is a breach", and $P$ the dyadic relation "can prevent", where $P(x,y)$ is read "$x$ can prevent $y$". Since it may not be immediately evident, how scope ambiguities relate to well-formed formulas in logic for strategic reasoning is discussed in the dataset section.

\begin{figure}[h]
\centering
\scriptsize
\begin{forest}
for tree={
    grow'=south,
    parent anchor=south,
    child anchor=north,
    align=center,
    l sep=1pt,
    s sep=1pt
}
[S
  [CoordP
    [S$_1$
      [NP [Bob]]
      [VP [has]]
    ]
    [Coord'
      [Conj [but]]
      [S$_2$
        [NP [Albert]]
        [VP [has not]]
      ]
    ]
  ]
  [{NP$_{\mathrm{shared}}$}
    [a strategy to ensure\\the machine is\\operative]
  ]
]
\end{forest}
\caption{Right Node Raising. Note that traces in VP are omitted for simplicity, same reasoning for the explosion of the complex NP "a strategy to ensure". Cartographic analysis provided in Appendix A.}
\end{figure}
\paragraph{Alternating-time Temporal Logic.}
Alternating-time Temporal Logic (ATL) and ATL$^*$ are strategic temporal logics
for reasoning about the abilities of coalitions of agents in multi-agent systems.
Both logics are interpreted over Concurrent Game Structures (CGS), the standard semantic models for logics for strategic reasoning. Formally, a CGS is a tuple
$G = (Agt, Ap, Act, St, s_0, d, \delta, \ell)$, where:
\begin{itemize}
  \item[-] $Agt$ is a finite, non-empty set of agents;
  \item[-] $Ap$ is a finite, non-empty set of atomic propositions;
  \item[-] $St$ is a non-empty set of states;
  \item[-] $s_0 \in St$ is the initial state;
  \item[-] $Act$ is a non-empty set of actions;
  \item[-] for each agent $a \in Agt$, $d_a : St \to 2^{Act} \setminus \{\emptyset\}$ gives the actions available to agent $a$;
  \item[-] $\delta : St \times \prod_{a \in Agt} Act \to St$ is the transition function;
  \item[-] $\ell : St \to 2^{Ap}$ is the labelling function.
\end{itemize}
Let $\lambda$ be a path, that is, a infinte sequence of states connected by successive transitions. We introduce the following notation:
\begin{itemize}
\renewcommand{\labelitemi}{-}
    \item $\lambda[i]$ denotes the state at position $i$ along $\lambda$, with indexing starting from $0$;
    \item $\lambda[i \dots j]$ denotes the subpath of $\lambda$ from position $i$ to position $j$;
    \item $\lambda[i \dots \infty]$ denotes the suffix of $\lambda$ starting at position $i$.
\end{itemize}
Additionally, we define a history $h \in St^+$ as a finite sequence of states that can occur in the system.
A perfect-recall strategy for agent $a$ is a function
$s_a : St^+ \to Act$ such that $s_a(h) \in d_a(\mathsf{last}(h))$ for every finite history $h$.
A collective strategy for a coalition $A \subseteq Agt$ is denoted by $s_A$.
The set of outcome paths starting from a state $q$ under $s_A$ is denoted by
$\mathsf{out}(q,s_A)$.

\paragraph{Syntax of ATL.}
The following syntax defines a well-formed formula in ATL:
\[
\varphi ::= p \mid \neg \varphi \mid \varphi \wedge \varphi
\mid \langle\!\langle A\rangle\!\rangle X \varphi
\mid \langle\!\langle A\rangle\!\rangle G \varphi
\mid \langle\!\langle A\rangle\!\rangle (\varphi \ U \ \varphi)
\]
Here, $p \in Ap$ and $A \subseteq Agt$. Additional boolean and temporal operators can be derived as usual. Temporal operator $F$ is classically derived.

\paragraph{Semantics of ATL.}
Given a CGS $M$ and a $q \in St$, the satisfaction relation $M,q \models \varphi$ is defined inductively as follows:
\begin{itemize}
  \item[-] $M,q \models p$ iff $p \in \ell(q)$;
  \item[-] $M,q \models \neg \varphi$ iff not $(M,q \models \varphi)$;
  \item[-] $M,q \models \varphi_1 \wedge \varphi_2$ iff
        $M,q \models \varphi_1$ and $M,q \models \varphi_2$;
  \item[-] $M,q \models \langle\!\langle A\rangle\!\rangle X \varphi$ iff
        $\exists s_A \ \forall \lambda \in \mathsf{out}(q,s_A):
        M,\lambda[1] \models \varphi$;
  \item[-] $M,q \models \langle\!\langle A\rangle\!\rangle G \varphi$ iff
        $\exists s_A \ \forall \lambda \in \mathsf{out}(q,s_A) \ \forall i \ge 0:
        M,\lambda[i] \models \varphi$;
  \item[-] $M,q \models \langle\!\langle A\rangle\!\rangle (\varphi_1 \ U \ \varphi_2)$ iff
        $\exists s_A \ \forall \lambda \in \mathsf{out}(q,s_A)$ such that
        $\exists i \ge 0$ with $M,\lambda[i] \models \varphi_2$ and
        $\forall j < i,\ M,\lambda[j] \models \varphi_1$.
\end{itemize}

\paragraph{Syntax of ATL$^*$.}
ATL$^\star$  distinguishes between state and path formulas, defined by:
\[
\begin{array}{l}
\varphi ::= p \mid \neg \varphi \mid \varphi \wedge \varphi
\mid \langle\!\langle A\rangle\!\rangle \gamma, \\
\gamma ::= \varphi \mid \neg \gamma \mid \gamma \wedge \gamma
\mid X\gamma \mid \gamma \ U \ \gamma .
\end{array}
\]
Here, $p \in AP$ and $A \subseteq Agt$. Temporal operators $G$ and $F$ are classically derived.

\paragraph{Semantics of ATL$^*$.}
Semantics of ATL$^*$ extends that for ATL as follows: 
%with the following clauses:
\begin{itemize}
  \item[-] $M,q \models \langle\!\langle A\rangle\!\rangle \gamma$ iff
        $\exists s_A \ \forall \lambda \in \mathsf{out}(q,s_A):
        M,\lambda \models \gamma$;
  \item[-] $M,\lambda \models \varphi$ iff $M,\lambda[0] \models \varphi$;
  \item[-] $M,\lambda \models X\gamma$ iff
        $M,\lambda[1 \dots \infty] \models \gamma$;
  \item[-] $M,\lambda \models \gamma_1 \ U \ \gamma_2$ iff
        $\exists i \ge 0$ s.t.
        $M,\lambda[i \dots \infty] \models \gamma_2$ and
        $\forall j < i,\ M,\lambda[j \dots \infty] \models \gamma_1$.
\end{itemize}




%\paragraph{Problem Statement.}